\section{소제 분석 (Ablation Study)}
\label{sec:소제분석}

제안 모델을 구성하는 핵심 기능 성분들의 기여도를 계량화하기 위해 소제 분석(Ablation Study)을 시행하였다. 오프라인 모방 모델, 지리적 residual prior, 위험 스위치 게이트 모듈, 그리고 Phase 13 P+의 상세 보정 기법의 효과를 차례로 검증하였다.

\subsection{주요 설계 모듈별 기여도 검증}
다음과 같은 기본 아블레이션 변종을 구성하여 PDR 및 전송 성능 변화를 비교하였다.
\begin{enumerate}
    \item \textbf{KD-only Student}: 지리적 residual prior를 제외하고, 오직 오프라인 교사 모델의 포워딩 결정 확률 분포만을 순수하게 KD로 모방한 MLP 모델.
    \item \textbf{KD + PPO Fine-tuning}: 오프라인 증류 완료 후, 로컬 보상만을 기반으로 1-hop 로컬 영역에서 PPO 알고리즘을 사용해 미세 튜닝(fine-tuning)을 거친 모델.
    \item \textbf{Geo-Residual KD (Phase 8)}: 지리적 진행 방향(Greedy progress)을 Prior 벡터로 고정하고, 학생 정책은 단지 토폴로지에 따른 보정값(Residual)만을 증류 학습하도록 유도한 모델.
    \item \textbf{Lite-GLOBE-P (No Switch)}: 상시적으로 인공신경망 기반의 예측(PRS) 차홉 결정을 수행하고 nominal 지리 결정을 생략한 모형.
    \item \textbf{Lite-GLOBE-P (\phaseTwelve)}: 지리 residual prior와 로컬 위험도 지표 기반의 스위칭 메커니즘을 모두 포함한 최종 모델.
\end{enumerate}

실험 결과, \textit{KD-only} 모델은 학습 단계와 기동 궤적이 다른 평가 시나리오에서 심각한 분포 변동(Distribution Shift) 오차를 보이며 PDR이 급격히 감소하였다. 로컬 무선 채널의 높은 변동성과 보상 희소성 때문에 \textit{PPO fine-tuning}은 뚜렷한 수렴 성능 향상을 제공하지 못하였다.

이에 반해, 지리 진행도를 우선 가이드하는 \textit{Geographic Residual} (Phase 8)을 결합한 순간 PDR이 KD-only 대비 약 12.4\% 증가하여, 물리적인 가이드를Prior로 주는 하이브리드 구조의 타당성이 입증되었다. 또한, 여기에 임계 위험 발생 시 대안 경로를 편향시키는 \textit{위험 스위치} 모듈(\phaseTwelve)을 융합함으로써, 급격한 노드 기동으로 인한 링크 이탈 확률을 사전에 방지하여 94\% 이상의 패킷 링크 손실 상황을 구제할 수 있었다.

\subsection{Phase 13 P+ 세부 모듈의 평가 방안}
Stochastic한 무선 환경에 적응하기 위해 추가 제안된 P+ 세부 모듈들의 상호 보완성을 검증하기 위해 다음의 4가지 ablation 시나리오를 설정하였다.
\begin{itemize}
    \item \textbf{No Link-Loss Gate}: 무선 페이딩에 따른 링크 손실 확률 $p_{i,\mathrm{keep}}$을 위험도 판별(Danger Score)에서 삭제한 모형.
    \item \textbf{No Energy-Aware Tie Breaking}: $\lambda_E = 0$으로 세팅하여, 안전도 유틸이 비슷할 때 근거리 에너지를 우선시하는tie-breaker를 비활성화한 모형.
    \item \textbf{No Drop Suppression}: 대안 경로가 존재해도 강제 전송을 조율하지 않고 드롭 판단에 따르도록 $\lambda_D = 0$을 가한 모형.
    \item \textbf{Top-1 Onward Only}: 상위-$k$ 평균 대신 단일 최선의 onward 예측 수명만을 사용한 모형.
\end{itemize}
이들 모듈의 최종 수치적 비교를 통해 각 P+ 메커니즘이 PDR 신뢰성 상승 및 에너지 절감에 미치는 기여가 확인될 예정이다.
